\documentclass[a4paper]{arrowhead}

\usepackage[yyyymmdd]{datetime}
\usepackage{etoolbox}
\usepackage[T1]{fontenc}
\usepackage[utf8]{inputenc}
\usepackage{multirow}
\usepackage{longtable}
\usepackage{amssymb}
\usepackage{float}
\usepackage{tikz}
\usetikzlibrary{positioning,arrows.meta,fit,calc,shapes.misc,shapes.geometric}

\renewcommand{\dateseparator}{-}
\setlength{\emergencystretch}{3em}
\lstset{basicstyle=\ttfamily\footnotesize,breaklines=true,columns=fullflexible,frame=single,rulecolor=\color{gray!50},backgroundcolor=\color{gray!6}}
\tikzset{
  ahsys/.style={draw=ArrowheadBlue,very thick,fill=ArrowheadBlue!6,rounded corners=2pt,minimum width=3.2cm,minimum height=1.1cm,align=center,font=\small\bfseries},
  ahext/.style={draw=ArrowheadGrey,thick,dashed,fill=gray!5,rounded corners=2pt,minimum width=3.2cm,minimum height=1.0cm,align=center,font=\small},
  ahsrv/.style={draw=ArrowheadPurple,thick,fill=ArrowheadPurple!8,rounded corners=6pt,minimum height=0.55cm,inner xsep=4pt,align=center,font=\scriptsize\ttfamily},
  ahcsm/.style={draw=ArrowheadTeal,thick,fill=ArrowheadTeal!8,rounded corners=6pt,minimum height=0.55cm,inner xsep=4pt,align=center,font=\scriptsize\ttfamily},
  ahdb/.style={draw=ArrowheadOrange,thick,fill=ArrowheadOrange!8,cylinder,shape border rotate=90,aspect=0.25,minimum height=0.9cm,align=center,font=\scriptsize},
  ahflow/.style={-{Stealth[length=2.2mm]},thick,ArrowheadGrey},
  ahlabel/.style={font=\scriptsize,fill=white,inner sep=1pt,align=center},
}

%% Special references
\newcommand{\fref}[1]{{\textcolor{ArrowheadBlue}{\hyperref[sec:functions:#1]{#1}}}}
\newcommand{\mref}[1]{{\textcolor{ArrowheadPurple}{\hyperref[sec:model:#1]{#1}}}}
\newcommand{\sref}[1]{{\textcolor{ArrowheadPurple}{#1}}}
\newrobustcmd\msubsection[2]{
  \addtocounter{subsection}{1}
  \addcontentsline{toc}{subsection}{\protect\numberline{\thesubsection}#1 \textcolor{ArrowheadPurple}{#2}}
  \subsection*{\thesubsection\quad#1 \textcolor{ArrowheadPurple}{#2}}
  \label{sec:model:#2}
}

\begin{document}

%% Arrowhead Document Properties
\ArrowheadTitle{WorkflowExecutor}

\ArrowheadType{System Description}
\ArrowheadTypeShort{SysD}
\ArrowheadVersion{5.4.0}
\ArrowheadDate{2026-11-23}
\ArrowheadAuthor{Documentation maintenance}
\ArrowheadStatus{REVIEW DRAFT}
\ArrowheadContact{Repository maintainers}
\ArrowheadFooter{\href{www.arrowhead.eu}{www.arrowhead.eu}}
\ArrowheadSetup
%%

%% Front Page
\begin{center}
  \vspace*{1cm}
  \huge{\arrowtitle}

  \vspace*{0.2cm}
  \LARGE{\arrowtype}
  \vspace*{1cm}

  
  \vspace*{\fill}

  

  \vspace*{1cm}
  \vspace*{\fill}

  % Front Page Abstract
  \begin{abstract}
    This document describes the WorkflowExecutor, an Eclipse Arrowhead support system that runs the steps of workflow sessions: it has providers orchestrated for itself, invokes them through equipment adapters and reports the outcome to the WorkflowManager. Documentation edition 5.4.0; system version 5.4.0 (framework), service versions 1.0.0.
  \end{abstract}

  \vspace*{1cm}
\end{center}

\newpage
%%

%% Table of Contents
\tableofcontents
\newpage
%%

\section{Overview}
\label{sec:overview}
The WorkflowExecutor (WfE) initiates the work. It receives work items from a WorkflowManager (push) or
claims them (pull), and runs each with itself as the consumer: lease handling, orchestration of the step's service
with the executor as requester, composition of the input from the session context, invocation of the provider through
an equipment adapter, and the report. Because the executor is the consumer, authorization, QoS filtering and provider
selection stay with the Arrowhead orchestrator and ConsumerAuthorization. An executor also implements the executor
interface of the Arrowhead Choreographer, so existing Choreographer plans can run on it.

\begin{figure}[ht!]
  \centering
  \resizebox{\textwidth}{!}{
\begin{tikzpicture}[node distance=4mm]
\node[ahsys,minimum height=6.75cm,minimum width=3.6cm] (sys) at (0,0) {WorkflowExecutor\\[-1pt]{\normalfont\scriptsize $\ll$system$\gg$}};
\node[ahsrv,anchor=west] (p0) at (3.4,3.00) {workflowExecution};
\draw[thick,ArrowheadPurple] (1.8,3.00) -- (p0.west);
\fill[ArrowheadPurple] (1.8,3.00) circle (1.4pt);
\node[ahsrv,anchor=west] (p1) at (3.4,2.25) {generalManagement};
\draw[thick,ArrowheadPurple] (1.8,2.25) -- (p1.west);
\fill[ArrowheadPurple] (1.8,2.25) circle (1.4pt);
\node[ahcsm,anchor=east] (c0) at (-3.4,3.00) {serviceDiscovery\\[-2pt]{\tiny\rmfamily ServiceRegistry}};
\draw[thick,ArrowheadTeal,-{Stealth[length=1.6mm]}] (c0.east) -- (-1.8,3.00);
\node[ahcsm,anchor=east] (c1) at (-3.4,2.25) {identity\\[-2pt]{\tiny\rmfamily Authentication}};
\draw[thick,ArrowheadTeal,-{Stealth[length=1.6mm]}] (c1.east) -- (-1.8,2.25);
\node[ahcsm,anchor=east] (c2) at (-3.4,1.50) {authorization\\[-2pt]{\tiny\rmfamily ConsumerAuthorization}};
\draw[thick,ArrowheadTeal,-{Stealth[length=1.6mm]}] (c2.east) -- (-1.8,1.50);
\node[ahcsm,anchor=east] (c3) at (-3.4,0.75) {serviceOrchestration\\[-2pt]{\tiny\rmfamily DynamicServiceOrchestration}};
\draw[thick,ArrowheadTeal,-{Stealth[length=1.6mm]}] (c3.east) -- (-1.8,0.75);
\node[ahcsm,anchor=east] (c4) at (-3.4,0.00) {workflowWorkItem\\[-2pt]{\tiny\rmfamily WorkflowManager}};
\draw[thick,ArrowheadTeal,-{Stealth[length=1.6mm]}] (c4.east) -- (-1.8,0.00);
\node[ahcsm,anchor=east] (c5) at (-3.4,-0.75) {workflowContext\\[-2pt]{\tiny\rmfamily WorkflowManager}};
\draw[thick,ArrowheadTeal,-{Stealth[length=1.6mm]}] (c5.east) -- (-1.8,-0.75);
\node[ahcsm,anchor=east] (c6) at (-3.4,-1.50) {workflowLease\\[-2pt]{\tiny\rmfamily WorkflowManager}};
\draw[thick,ArrowheadTeal,-{Stealth[length=1.6mm]}] (c6.east) -- (-1.8,-1.50);
\node[ahcsm,anchor=east] (c7) at (-3.4,-2.25) {(step services)\\[-2pt]{\tiny\rmfamily Equipment / application systems}};
\draw[thick,ArrowheadTeal,-{Stealth[length=1.6mm]}] (c7.east) -- (-1.8,-2.25);
\node[ahcsm,anchor=east] (c8) at (-3.4,-3.00) {choreographerExecutor\\[-2pt]{\tiny\rmfamily Choreographer}};
\draw[thick,ArrowheadTeal,-{Stealth[length=1.6mm]}] (c8.east) -- (-1.8,-3.00);
\node[font=\scriptsize\bfseries,ArrowheadPurple] at (5.0,3.90) {produced};
\node[font=\scriptsize\bfseries,ArrowheadTeal] at (-5.0,3.90) {consumed};
\end{tikzpicture}
  }
  \caption{SysML block view of the WorkflowExecutor: produced services (right) and consumed services with their providers (left).}
  \label{fig:block}
\end{figure}

The rest of this document is organized as follows.
In Section~\ref{sec:prior_art}, we reference major prior art capabilities of the system.
In Section~\ref{sec:use}, we describe the intended usage of the system.
In Section~\ref{sec:properties}, we describe fundamental properties provided by the system.
In Section~\ref{sec:delimitations}, we describe delimitations of capabilities of the system.
In Section~\ref{sec:services}, we describe the abstract service functions consumed or produced by the system.
In Section~\ref{sec:security}, we describe the security capabilities of the system.

\newpage
\subsection{Significant Prior Art}
\label{sec:prior_art}
Derhamy's workflow executor, customised per workstation and talking to nothing outside its work
centre~\cite{derhamy2018workflow}; the Choreographer executor of Arrowhead 4.6 and 5.x; Varga's execution primitive
$compile \rightarrow O \rightarrow consume$ with the step's consumer as requester~\cite{varga2026}.

\subsection{How This System Is Meant to Be Used}
\label{sec:use}
One or more executors are deployed per workstation. Each declares its capabilities (e.g. \texttt{heating},
\texttt{nutrunner}), the services it can run (\texttt{*} for any) and its workstation; these are registered as
metadata of \texttt{workflowExecution}. In push mode the WorkflowManager selects the executor; in pull mode the
executor polls \texttt{workflowWorkItem.claim}. With \texttt{workflow.choreographer.enabled} the executor also
registers with the Choreographer (\texttt{choreographerExecutor.register}) and runs Choreographer plan steps.

\subsection{System functionalities and properties}
\label{sec:properties}
\subsubsection{Functional properties of the system}
\begin{itemize}
\item Run a work item: lease reserve/poll or renew, orchestration (pull, matchmaking, inter-cloud only when the step
allows it) at the configured orchestrator, input composition $\kappa_s(\rho, \xi|dep(s))$, invocation, lease release,
report (BOUND, BIND\_FAILED, SUCCESS, ERROR, FATAL\_ERROR, ABORTED).
\item Equipment adapter SPI with the generic HTTP/HTTPS JSON adapter (\texttt{generic\_http}, \texttt{generic\_https}).
\item Pull loop with configurable polling interval; push endpoint; abort and status of running work items.
\item Choreographer executor compatibility: \texttt{service-info}, \texttt{start}, \texttt{abort}, reporting with
\texttt{notify-step-done}; registration and unregistration with the Choreographer.
\end{itemize}

\subsubsection{Configuration of system properties}
The system is configured with Spring properties; the main ones are listed in the SysDD. The workflow-specific behaviour
is configured under the prefix \texttt{workflow.}; the security profile (TLS, authentication, authorization) is the
common Arrowhead~5 profile.

\subsubsection{Data stored by the system}
No workflow data is persisted. Work items in progress are kept in memory; the common-library log table is kept in an
embedded H2 database.

\subsubsection{Non functional properties}
\begin{itemize}
\item \emph{Security:} see Section~\ref{sec:security}. Providers see the executor's own certificate and identity.
\item \emph{Latency:} one orchestration and one invocation per step, plus two calls to the WorkflowManager (context and
report); work items run on a thread pool (\texttt{workflow.executor.threads}, default 8).
\item \emph{Energy and power saving:} in pull mode the executor polls (default every 500\,ms); push mode is idle
between work items.
\end{itemize}

\subsubsection{Stateful or stateless}
Stateless with respect to workflows: the authoritative state is in the WorkflowManager. After a restart
unknown work items are reported as \texttt{UNKNOWN} by \texttt{get-status}, and the WorkflowManager's timeouts apply.

\subsection{Important Delimitations}
\label{sec:delimitations}
\begin{itemize}
\item Only HTTP/HTTPS JSON equipment; MQTT, CoAP and OPC UA adapters are not implemented (the SPI exists).
\item TOKEN\_AUTH provider interfaces (orchestration tokens) are not used.
\item Never consumes services outside its local cloud unless the step's policy sets \texttt{localCloudOnly = false}.
\item Does not store recipes or results; does not talk to products.
\end{itemize}

\newpage
\section{Services}
\label{sec:services}
This section lists the produced and consumed services. Each produced service is specified by an SD and an HTTP/TLS/JSON
IDD in this documentation set.

\subsection{Produced services}

{\small
\begin{longtable}{|p{3.8cm}|p{6.2cm}|p{5.8cm}|}

\hline
\rowcolor{gray!33} Service & Operations & Documents \\ \hline
\endhead
\texttt{workflowExecution} & service-info, start, abort, get-status & \texttt{WorkflowExecutor-v5.4.0-workflowExecution-SD}, \texttt{-HTTP-IDD} \\ \hline
\texttt{generalManagement} & get-log, get-config & \texttt{Workflow-v5.4.0-generalManagement-SD}, \texttt{-HTTP-IDD} \\ \hline
\texttt{Choreographer executor interface} & service-info, start, abort (not registered; under \textbackslash{}texttt\{/workflowexecutor/choreographer\}) & Choreographer executor IDD (Choreographer 5.x) \\ \hline
\end{longtable}
}


\subsection{Consumed services}
All three workflow systems use the core systems the same way: at start-up they log in at the Authentication
system (\texttt{identity}) and register their services in the ServiceRegistry (\texttt{serviceDiscovery.register});
on every incoming request they verify the caller's identity token (\texttt{identity.identity-verify}) and, with
\texttt{service.authorization.enabled}, the caller's grant (\texttt{authorization.verify} of ConsumerAuthorization);
at shutdown they revoke their services and log out.

{\small
\begin{longtable}{|p{3.4cm}|p{3.2cm}|p{3.6cm}|p{5.6cm}|}

\hline
\rowcolor{gray!33} Service & Provider & Operations & Purpose \\ \hline
\endhead
\texttt{serviceDiscovery} & ServiceRegistry & register, revoke & Registration of \texttt{workflowExecution} with capability metadata. \\ \hline
\texttt{identity} & Authentication & login, logout, identity-verify & Own identity token; verification of callers. \\ \hline
\texttt{authorization} & ConsumerAuthorization & verify & Provider-side authorization. \\ \hline
\texttt{serviceOrchestration} & DynamicServiceOrchestration & pull & Provider of each step, with the executor as requester. \\ \hline
\texttt{workflowWorkItem} & WorkflowManager & claim, report & Pull and reports. \\ \hline
\texttt{workflowContext} & WorkflowManager & get-parameters, get-results, put-result & $\rho$, $\xi|dep(s)$, $\xi(s)$. \\ \hline
\texttt{workflowLease} & WorkflowManager & reserve, poll, renew, release & Shared resources. \\ \hline
(step services) & Equipment / application systems & as orchestrated & The work itself. \\ \hline
\texttt{choreographerExecutor} & Choreographer & register, unregister, notify-step-done & Only with \texttt{workflow.choreographer.enabled}. \\ \hline
\end{longtable}
}


\newpage
\section{Security}
\label{sec:security}
The WorkflowExecutor runs in Arrowhead secure mode in every cloud deployment; an insecure mode exists only for development and
unit tests. It handles only Arrowhead-compliant X.509 certificates of its local cloud; others are refused during the
TLS handshake.

\subsection{Security model}
\begin{itemize}
\item \textbf{Secure mode only in deployment.} The deployment configuration (\texttt{deployment/*.properties}) enables
TLS with client authentication (\texttt{server.ssl.client-auth=need}, TLS 1.3 and 1.2), outsourced authentication with
identity tokens bound to the client certificate, interface policy \texttt{CERT\_AUTH} and the CRL check. The packaged
defaults without TLS exist for development and unit tests (\texttt{workflow.standalone}) and must not be used in a
cloud.
\item \textbf{Certificates.} The system has an Arrowhead system certificate
\texttt{<System>.<cloud>.<organisation>.arrowhead.eu} issued by the local cloud CA, kept in a PKCS\#12 key store; the
trust store holds the cloud CA. Non-Arrowhead certificates are refused at the TLS layer.
\item \textbf{Provider-side authorization.} Every operation of the registered services (other than management) is
checked against ConsumerAuthorization (\texttt{authorization.verify}) for the calling system, the service definition
and the operation (route tables \texttt{ServiceOperations} and \texttt{ConfiguratorOperations}). A missing grant gives
403.
\item \textbf{Management.} \texttt{generalManagement} and \texttt{workflowManagement} are protected by
\texttt{management.policy=sysop-only}.
\item \textbf{Data protection.} All service traffic is encrypted by TLS. The WorkflowManager's database connection uses
TLS with server certificate verification (\texttt{sslMode=verify-full}). Secrets (store, identity and database
passwords, the console password hash) come from container secrets (\texttt{configtree:/run/secrets/}), not from the
image.
\end{itemize}


\subsection{Object-level checks}
Only the configured WorkflowManager (or a sysop) may start or abort work items, and only the configured Choreographer
may use the Choreographer executor interface. Providers of step services see the executor's own certificate and must
grant the executor, not the WorkflowManager.

\newpage
\bibliographystyle{IEEEtran}
\bibliography{bibliography}

\newpage

\section{Revision History}
\subsection{Amendments}

\noindent\begin{tabularx}{\textwidth}{| p{1cm} | p{2.2cm} | p{1.6cm} | X | p{3.6cm} |} \hline
\rowcolor{gray!33} No. & Date & Version & Subject of Amendments & Author \\ \hline
1 & 2026-11-23 & \arrowversion & Initial edition: system description of the WorkflowExecutor, written from the Java reference implementation (ah-Workflow-java). & Documentation maintenance \\ \hline
\end{tabularx}

\subsection{Quality Assurance}

\noindent\begin{tabularx}{\textwidth}{| p{1cm} | p{2.2cm} | p{1.6cm} | X |} \hline
\rowcolor{gray!33} No. & Date & Version & Approved by \\ \hline
1 & & \arrowversion & Pending review \\ \hline
\end{tabularx}

\end{document}
